% id: 135-294
% book: 개념원리 확률과 통계 · 14 이항분포
% section: 개념원리 익히기
% figure: none
% answer: ⑴ $\mathrm{B}\!\left(10{,}\mkern3mu\mathopen{}\dfrac{1}{2}\right)$ ⑵ 이항분포를 따르지 않는다. ⑶ $\mathrm{B}(50{,}\mkern3mu\mathopen{}0.85)$ ⑷ $\mathrm{B}(100{,}\mkern3mu\mathopen{}0.4)$
% answer_source: 답지
% variant_level: 0 (원본 전사)
% 이 파일은 build.mjs 가 items.json 에서 만든다. 고칠 때는 items.json 을 고친다.
\smash{\raisebox{1.5mm}{\dmkichul{개념원리 확률과 통계 135쪽 294번}}}
다음 확률변수 $X$가 이항분포를 따르는지 확인하고, 이항분포를 따르면 $X$의 확률분포를 $\mathrm{B}(n{,}\mkern3mu\mathopen{}p)$ 꼴로 나타내시오.
\dmsubprob{1} $10$개의 주사위를 동시에 던질 때, 짝수의 눈이 나오는 주사위의 개수 $X$
\dmsubprob{2} $3$개의 당첨 제비를 포함한 $20$개의 제비 중에서 임의로 $2$개를 차례대로 뽑을 때, 뽑은 당첨 제비의 개수 $X$ \cond{꺼낸 제비는 다시 넣지 않는다.}
\dmsubprob{3} 자유투 성공률이 $0.85$인 어느 농구 선수가 자유투를 $50$번 던져서 성공하는 횟수 $X$
\dmsubprob{4} 재구매율이 $40\mkern3mu \%$인 화장품을 $100$명에게 판매하였을 때, 재구매하는 인원수 $X$
